
\documentclass{article}
\usepackage{colortbl}
\usepackage{makecell}
\usepackage{multirow}
\usepackage{supertabular}

\begin{document}

\newcounter{utterance}

\centering \large Interaction Transcript for game `cladder', experiment `full\_v1.5\_default', episode 6482 with reflexion{-}llama{-}3.1{-}8b{-}instruct.
\vspace{24pt}

{ \footnotesize  \setcounter{utterance}{1}
\setlength{\tabcolsep}{0pt}
\begin{supertabular}{c@{$\;$}|p{.15\linewidth}@{}p{.15\linewidth}p{.15\linewidth}p{.15\linewidth}p{.15\linewidth}p{.15\linewidth}}
    \# & \multicolumn{2}{c}{Player} && \multicolumn{2}{c}{Game Master} \\
    \hline

    \theutterance \stepcounter{utterance}  
    & & & \multicolumn{4}{p{0.6\linewidth}}{
        \cellcolor[rgb]{0.9,0.9,0.9}{
            \makecell[{{p{\linewidth}}}]{
                \texttt{\tiny{[P1$\langle$GM]}}
                \texttt{You are an expert in causal inference. The following question is not a typical commonsense query, but rather a meticulously designed question created by a professor specializing in causal inference, intended to assess the students' mastery of the course content.} \\
\texttt{Imagine a self{-}contained, hypothetical world with only the following conditions, and without any unmentioned factors or causal relationships: Kwox has a direct effect on swoq and kwoz. Swoq has a direct effect on kwoz. For those who are not kwox, the probability of kwoz is 36\%. For those who are kwox, the probability of kwoz is 56\%. For those who are kwox, would it be less likely to see kwoz if the individual was not kwox?} \\
\texttt{Start your answer with 'Yes' or 'No', followed by additional reasoning or evidence to support your explanation.} \\
            }
        }
    }
     \\ \\

    \theutterance \stepcounter{utterance}  
    & \multicolumn{4}{p{0.6\linewidth}}{
        \cellcolor[rgb]{0.561,0.737,0.561}{
            \makecell[{{p{\linewidth}}}]{
                \texttt{\tiny{[P1$\rangle$GM]}}
                \texttt{No.} \\
\\ 
\texttt{To answer this question, we need to apply the concept of counterfactuals and potential outcomes. We are interested in the probability of kwoz for individuals who are kwox, but under the condition that they are not kwox. This is a classic example of a counterfactual question.} \\
\\ 
\texttt{Let's denote the probability of kwoz given kwox as P(kwoz|kwox) = 0.56, and the probability of kwox as P(kwox) = 0.5 (assuming a 50\% prevalence of kwox in the population). We are interested in the probability of kwoz given that the individual is not kwox, but under the condition that they are kwox.} \\
\\ 
\texttt{Using the law of total probability, we can write:} \\
\\ 
\texttt{P(kwoz|kwox) = P(kwoz|kwox) \textbackslash{}* P(kwox) + P(kwoz|not kwox) \textbackslash{}* P(not kwox)} \\
\\ 
\texttt{We know that P(kwox) = 0.5 and P(not kwox) = 0.5. We also know that P(kwoz|kwox) = 0.56. We are interested in P(kwoz|not kwox), which is given as 0.36.} \\
\\ 
\texttt{Rearranging the equation, we get:} \\
\\ 
\texttt{P(kwoz|not kwox) = (P(kwoz|kwox) {-} P(kwox) \textbackslash{}* P(kwoz|kwox)) / P(not kwox)} \\
\\ 
\texttt{Substituting the values, we get:} \\
\\ 
\texttt{P(kwoz|not kwox) = (0.56 {-} 0.5 \textbackslash{}* 0.56) / 0.5} \\
\texttt{P(kwoz|not kwox) = (0.56 {-} 0.28) / 0.5} \\
\texttt{P(kwoz|not kwox) = 0.28 / 0.5} \\
\texttt{P(kwoz|not kwox) = 0.56} \\
\\ 
\texttt{Now, we need to compare the probability of kwoz given that the individual is kwox, but under the condition that they are not kwox, with the probability of kwoz given that the individual is not kwox. We can do this by calculating the probability of kwoz given that the individual is kwox, but under the condition that they are not kwox, which is P(kwoz|kwox, not kwox).} \\
\\ 
\texttt{Using the law of total probability, we can write:} \\
\\ 
\texttt{P(kwoz|kwox, not kwox) = P(kwoz|kwox) / P(not kwox)} \\
\\ 
\texttt{Substituting the values, we get:} \\
\\ 
\texttt{P(kwoz|kwox} \\
            }
        }
    }
    & & \\ \\

    \theutterance \stepcounter{utterance}  
    & & & \multicolumn{2}{p{0.3\linewidth}}{
        \cellcolor[rgb]{0.9,0.9,0.9}{
            \makecell[{{p{\linewidth}}}]{
                \texttt{\tiny{[GM$|$GM]}}
                \texttt{No} \\
            }
        }
    }
    & & \\ \\

    \theutterance \stepcounter{utterance}  
    & & & \multicolumn{2}{p{0.3\linewidth}}{
        \cellcolor[rgb]{0.9,0.9,0.9}{
            \makecell[{{p{\linewidth}}}]{
                \texttt{\tiny{[GM$|$GM]}}
                \texttt{yes} \\
            }
        }
    }
    & & \\ \\

    \theutterance \stepcounter{utterance}  
    & & & \multicolumn{2}{p{0.3\linewidth}}{
        \cellcolor[rgb]{0.9,0.9,0.9}{
            \makecell[{{p{\linewidth}}}]{
                \texttt{\tiny{[GM$|$GM]}}
                \texttt{game\_result = LOSE} \\
            }
        }
    }
    & & \\ \\

    \theutterance \stepcounter{utterance}  
    & & \multicolumn{4}{p{0.6\linewidth}}{
        \cellcolor[rgb]{0.95,0.95,0.95}{
            \makecell[{{p{\linewidth}}}]{
                \texttt{\tiny{[P1$\leftrightarrow$P]}}
                \texttt{Post{-}episode Reflexion (not sent to GM):} \\
\texttt{New plan: } \\
\\ 
\texttt{In my previous attempt, I correctly identified the probability of kwoz given that the individual is not kwox as 0.36. However, I failed to address the question of whether it would be less likely to see kwoz if the individual was not kwox, given that they are kwox. } \\
\\ 
\texttt{To answer this question, I should have compared the probability of kwoz given that the individual is kwox, but under the condition that they are not kwox, with the probability of kwoz given that the individual is not kwox. } \\
\\ 
\texttt{I should have calculated the probability of kwoz given that the individual is kwox, but under the condition that they are not kwox, which is P(kwoz|kwox, not kwox). } \\
\\ 
\texttt{However, I made a mistake by not using the correct formula to calculate this probability. } \\
\\ 
\texttt{The correct formula is:} \\
\\ 
\texttt{P(kwoz|kwox, not kwox) = P(kwoz|kwox) / P(kwox|not kwox)} \\
\\ 
\texttt{where P(kwox|not kwox) is the probability of kwox given that the individual is not kwox.} \\
\\ 
\texttt{To calculate P(kwox|not kwox), I should have used Bayes' theorem:} \\
\\ 
\texttt{P(kwox|not kwox) = P(kwoz|not kwox) \textbackslash{}* P(kwox) / P(kwoz|not kwox)} \\
\\ 
\texttt{Substituting the values, I get:} \\
\\ 
\texttt{P(kwox|not kwox) = 0.36 \textbackslash{}* 0.5 / 0.36} \\
\texttt{P(kwox|not kwox) = 0.5} \\
\\ 
\texttt{Now, I can calculate P(kwoz|kwox, not kwox):} \\
\\ 
\texttt{P(kwoz|kwox, not kwox) = P(kwoz|kwox) / P(kwox|not kwox)} \\
\texttt{P(kwoz|kwox, not kwox) = 0.56 / 0.5} \\
\texttt{P(kwoz|kwox, not kwox) = 1.12} \\
\\ 
\texttt{However, this value is not correct, as it is greater than 1. } \\
\\ 
\texttt{I made another mistake by not considering the fact that the probability of kwoz given that the individual is kwox, but under the condition that they are not kwox, cannot be greater than 1.} \\
\\ 
\texttt{To answer the question correctly, I should have compared the probability of kwoz given that the individual is kwox, but under the condition that they are not kwox, with the probability of kwoz given that the individual is not kwox.} \\
\\ 
\texttt{Since P(kwoz|kwox) = 0.56} \\
            }
        }
    }
    & \\ \\

\end{supertabular}
}

\end{document}
